%% kapa: exam=1 round=21 year=2010 period=1 form=A subject=civil_law subject_ko=민법 seq=22 number=22
%% source: Q-Net 감정평가사 기출문제 (cst00309) artlSeq=5024181 file=extracted/1008/5024181/2010년 제21회 감정평가사 1차시험 기출문제/1교시.hwp
%% conversion: kordoc 4.24.51 -> tools/gen1.py (auto); review=proofread-sonnet; flags=none
\begin{question}{22}
주위토지통행권에 관한 설명으로 옳지 않은 것은? (다툼이 있으면 판례에 의함)
\begin{choices}
  \item 토지소유자가 주위의 토지를 통행 또는 통로로 하지 않으면 공로에 출입할 수 없는 경우는 물론 과다한 비용을 요할 때에도 주위토지통행권이 인정된다.
  \item 주위토지통행권의 범위는 현재 뿐만 아니라 장래의 이용 상황까지 미리 대비하여 결정할 수 있음이 원칙이다.
  \item 이미 기존의 통로가 있더라도 그것이 당해 토지의 이용에 부적합하여 실제로 통로로서의 충분한 기능을 하지 못하고 있는 경우에도 주위토지통행권이 인정된다.
  \item 통행권자가 통행함에 그치지 않고 통행지를 배타적으로 점유하고 있다면 통행지소유자는 통행권자에 대하여 그 인도를 청구할 수 있다.
  \item 주위토지통행권이 발생하였더라도 후에 그 토지에 접하는 공로가 개설됨으로써 주위토지통행권을 인정할 필요성이 없어진 때에는 그 통행권은 소멸한다.
\end{choices}
\end{question}
\answer{2}
